% app-revisions.tex -- the release history, and the rule for keeping it
\chapter{Revision history}
\label{ch:revisions}

This appendix records every release: what changed, and --- for anything that
moved a number --- how much it moved, on which case. That second half is the
part with teeth. An entry that says a result changed without saying by how much
is not an entry, because it leaves you unable to decide whether your own work
is affected; \cref{sec:policy:results} is the commitment this list exists to
keep. Releases are listed newest first. \textbf{Version 1.0 is the first
public release}, and the numbering starts here. The same entries are kept in
\file{CHANGELOG.md} at the top of the repository, whose section for a version
is the text of that release on GitHub.

\section{Version 1.0}
\label{sec:rev:1-0}

\textbf{Release date:} the date of the \code{v1.0.0} tag. This manual is built
from the tag and attached to that release, so the release it documents is the
one it came from; the repository's releases page carries the date.

The first public release. What follows describes what it contains.

\begin{note}[Two names an input file and a reader must match]
  An input file must read \key{SolverName}~$=$~\val{"Yuzen-1.0"}; anything
  else is rejected with a range error rather than silently accepted. And the
  VTK output opens with the element \code{<YuzenData ...>}, which is what a
  script parsing those files should match on.
\end{note}

\subsection*{Solver}

\begin{itemize}
  \item \textbf{Three models are implemented} and are selected by
    \key{Physics.ConcField.Model} and \key{Grid.Dimension} through the factory
    \code{MakeSolverModel}: well-mixed in one dimension, well-mixed in two
    dimensions, and \val{"CZ"}, which resolves the concentration through the
    thickness (\cref{ch:concentration}).
  \item Six flux contributions --- capillary pressure, disjoining pressure,
    vertical gravity, in-plane gravity, centrifugal force and solutal
    capillarity --- together with evaporation and a moving source
    (\cref{ch:terms}). Each is switched on and off independently, and each
    dimensionless group is exactly zero when its term is off.
  \item Both evolution equations are advanced with Crank--Nicolson. The
    thickness equation is driven to the implicit solution by a
    Newton--Kantorovich outer iteration whose convergence test is taken on the
    larger of the thickness and concentration corrections, so the converged
    state is the fully implicit coupled solution of the two
    (\cref{sec:time:coupling}). The time step is optionally adaptive.
  \item The convection schemes are chosen in the input
    (\key{Numerics.Convection}, \cref{sec:ref-numerics:convection}): the
    Koren-limited scheme is the default for the upwinded fluxes of the film
    height, which keeps the face height between its neighbors, and UTOPIA
    for the concentration; KK, van Leer, fourth-order central and others can
    be selected (\cref{sec:verification:schemes}).
  \item Two interchangeable linear-solver backends sit behind
    \code{ILinearSolverBackend} (\cref{ch:linearsolver}).
  \item \textbf{Thread invariance is asserted by the test suite}, not merely
    claimed: \code{ctest} runs the same binary in three separate processes, at
    the environment's own thread count, at one thread and at four, and the
    reductions and the linear solver are compared with \code{==}
    (\cref{sec:verification:linalg}). The \val{"amgcl"} backend is the
    documented exception (\cref{sec:linsolve:choosing}).
  \item \textbf{More threads do not make a run slower within a socket}: a GCC
    build links LLVM's \code{libomp} when it is installed, and linear solves
    take threads by their size (\env{YUZEN\_GRAIN}); a run uses at most the
    cores of one socket (\cref{sec:parallel:runtime,sec:parallel:cap}).
  \item \textbf{The disjoining-pressure calibration is verified in closed
    form}: over 36 parameter combinations the coefficient $\Dpn$ agrees with
    the analytic expression to within $10^{-14}$ relative, and the recovered
    equilibrium contact angle to within $10^{-12}$
    (\cref{sec:verification:disjoining}). The discretization of capillary plus
    disjoining pressure converges at order 1.94 on the analytic equilibrium
    droplet.
\end{itemize}

\subsection*{Input and output}

\begin{itemize}
  \item Input is a TOML file (\cref{ch:inputformat}). A block-structured \file{.inp}
    file is also read and produces the same tree, but its format is not
    specified here (\cref{sec:inputformat:legacy}). \key{SolverName} must read
    \val{"Yuzen-1.0"}; the accepted value names \code{MAJOR.MINOR} only
    (\cref{sec:policy:scheme}).
  \item A result directory carries a verbatim copy of the input file,
    \file{JobInfo.txt}, \file{Warnings.txt}, \file{FatalErrors.txt}, the field
    output in VTK and tab-separated text, \file{History.dat},
    \file{CalcInfo.dat} and any probe files (\cref{sec:output:dir}). The run
    log itself goes to standard output and to no file
    (\cref{sec:output:running}). \file{field*.dat} is also the ASCII restart
    format, and \file{History.dat} and \file{CalcInfo.dat} open with a header
    line naming their columns.
  \item The field output is written every \key{TimeInterval} or at listed
    times (\key{Times}, or a \key{LogSchedule} spaced logarithmically), each
    of them capping only the step that would pass it, and with
    \key{ExactTimes} landed on (\cref{sec:ref-output:listed}).
  \item \key{WritePotential} adds the driving potential $\varphi$ to the
    well-mixed field output (\cref{sec:output:potential}).
  \item The capillary pressure, the disjoining pressure and the vertical
    gravity share one positive face mobility, so the dissipation of the
    scheme is a sum of squares and a steady film the minimum of the discrete
    free energy (\cref{sec:disc:positivity-one}).
  \item \key{Output.Energy} writes \file{Energy.dat}: the free energy of the
    film, the dissipation of the scheme and their balance, second order in
    $\Delta t$ (\cref{sec:output:energy}).
  \item \key{Physics.Slip} gives the substrate a Navier slip length in the
    well-mixed model: mobility $h^3/3 + \lambda h^2$, Marangoni
    $h^2/2 + \lambda h$, a positive face mobility built for it
    (\cref{sec:ref-physics:slip,sec:disc:positivity-slip}). Without the
    block every number is as it was (\code{compare.py} 0 on every shipped
    case).
  \item An entry the parser does not read --- a misspelled key, a key in the
    wrong table, a misspelled table --- stops the run with its full path and
    the likely intent (\cref{sec:inputformat:unknown}).
    \code{yuzen --check-input} makes that check and the others on a file
    without running it.
  \item An optional block left out of the file reads as an empty one; a
    missing \key{Physics} sub-block is a term switched off, and a missing
    \key{Physics.ConcField} is \key{Model}~$=$~\val{"None"}
    (\cref{sec:inputformat:defaults}). \key{ExponentN} and
    \key{ExponentM} of \key{Physics.DisjoiningPressure} default to
    $(3,2)$ inside a block that is written as well.
    \key{Grid.Xmax}~$>$~\key{Grid.Xmin}, \key{ChangeFactor} in $(0:1)$
    and \key{KeepSteps}, the \key{MaxIterations} of both linear solvers
    and \key{SubRegions} of at least 1 are checked as the manual states.
  \item \key{Numerics.MaxSteps} stops the run; a run that reaches it says
    \code{Step count reached Numerics.MaxSteps.} rather than
    \code{Time reached to the EndTime.}
    (\cref{sec:ref-numerics:numerics}). It defaults to \val{0}, no limit, and
    is counted in 64 bits, so any step count a run can reach is representable.
  \item \key{Grid.Nx}, \key{Ny} and \key{Nz} accept up to one million cells
    each, with the total
    $(\mathrm{Nx}+4)(\mathrm{Ny}+4)(\mathrm{Nz}+4)$ capped at $5\times10^{7}$
    so that the grid's linear index stays inside an \code{int}; a 1-D run can
    therefore be made large enough to occupy a large machine
    (\cref{sec:parallel:kit}).
  \item \file{JobInfo.txt} carries a \code{\#\#\#\#\#\# Performance
    \#\#\#\#\#\#} section: the wall-clock time of the time-step loop alone,
    the number of steps taken, and the total Newton and linear iterations.
    \code{Elapsed Time} covers the whole process, including grid
    construction and the initial output, which makes it useless as a scaling
    measurement on a large grid.
\end{itemize}

\subsection*{Documentation}

\begin{itemize}
  \item This manual, the authoritative description of the solver, with the
    shorter companion documents \file{docs/theory.md}, \file{docs/numerics.md}
    and \file{docs/validation.md} (\cref{sec:extending:submitting}).
  \item \Cref{ch:verification} records the evidence for every claim
    \emph{together with how strong that evidence is}, and every number in it
    comes with the command that reproduces it.
  \item \file{tools/scaling/} is a shared-memory scaling kit --- generated
    cases that fix the work per run, a restartable sweep script, and a
    summarizer (\cref{sec:parallel:kit}). It exists because the scaling
    numbers in \cref{tab:scaling} were measured on one machine and cannot
    speak for another; the kit is how a reader gets numbers that speak for
    theirs.
\end{itemize}

\subsection*{Known gaps}

\begin{itemize}
  \item \Cref{ch:notverified} is the complete list of what has not been
    checked, and it is part of the release rather than an apology for it. The
    largest items are that the contact-line model is verified in one dimension
    only, that the order in time is measured with every term on but not for
    every combination of switches, that the
    \val{"amgcl"} backend --- although its agreement with \val{"eigen"} and its
    thread-to-thread difference are measured
    (\cref{sec:verification:backends,sec:verification:amgclthreads}) --- is
    exercised by no test in the suite, and that \val{"FreshInlet"} and
    \val{"Outlet"} are checked against each other and against the balances
    through them, not against a known inflow.
\end{itemize}

\section{Keeping this list}
\label{sec:rev:keeping}

For the maintainer. An entry is not free-form. It must carry, at minimum:

\begin{itemize}
  \item the \textbf{version}, in full, \code{MAJOR.MINOR.PATCH};
  \item the \textbf{release date};
  \item \textbf{one line per user-visible change} --- a change nobody outside
    the repository can observe does not belong here;
  \item and for any change that altered numerical output, \textbf{the case it
    was measured on and the size of the difference}. Quantified, not
    characterized: ``\file{allphys\_1d\_reg}, \code{maxH} moved by
    $3.2\times10^{-4}$ relative'' is an entry, and ``results change slightly'' is
    not.
\end{itemize}

Prefix each line with what kind of change it is --- \code{fix}, \code{add},
\code{change}, \code{doc} --- and say in the same line whether input files are
affected. A renamed or removed key must be named here with both spellings, so
that a user whose input file is stopped on the old one finds the new one; a
changed default must be named with both values, because nothing in the parser
can tell a user about that (\cref{sec:policy:compat}).

The format is then unambiguous:

\begin{plain}
1.0.1  2026-10-03
  fix: the evaporation cap was applied after the source term instead of before
       it, so a cell under the nozzle could lose more than half its thickness
       in one step.  allphys_1d_reg: History.dat Volume 1.1e-5, field maxH
       3.2e-4 (both relative).  Input files unaffected.
  doc: numerics.md, the section on the evaporation cap.
\end{plain}

A release in which nothing moved says so explicitly, in the same place and just
as plainly: \code{no case changed; compare.py reads 0.000e+00 throughout}.
That sentence is worth as much to a reader as any of the others, and it is the
one most often left out.
